% concatenated from Searching_for_Courter11.tex
\documentclass[11pt]{article}
\usepackage{amsmath,amssymb,amsthm}
\DeclareMathOperator{\End}{End}
\DeclareMathOperator{\Hom}{Hom}
\DeclareMathOperator{\Soc}{Soc}
\DeclareMathOperator{\rad}{rad}
\usepackage{geometry}
\geometry{margin=1in}
\usepackage{graphicx}
\usepackage{float}
\usepackage{caption}
\usepackage{hyperref}
\usepackage[T1]{fontenc}
\usepackage[utf8]{inputenc}
\theoremstyle{remark}
\usepackage{lineno}
\newtheorem*{example}{Example}
\newtheorem*{question}{Question}
\usepackage{array}
\usepackage{tabularx}
\theoremstyle{plain}
\newtheorem{theorem}{Theorem}
\newtheorem{proposition}[theorem]{Proposition}
\newtheorem{lemma}[theorem]{Lemma}
\newtheorem{corollary}[theorem]{Corollary}
\usepackage{listings}
\usepackage{xcolor}

\lstset{
	basicstyle=\ttfamily\small,
	breaklines=true,
	columns=fullflexible,
	keepspaces=true,
	frame=single
}

\theoremstyle{definition}
\newtheorem{definition}[theorem]{Definition}

\newtheoremstyle{myremark}
{6pt}
{6pt}
{\normalfont}
{}
{\bfseries}
{.}
{0.5em}
{}

\theoremstyle{myremark}
\newtheorem{remark}{Remark}


%\newtheorem{remark}{Remark}
%\usepackage[colorlinks=true, linkcolor=blue, citecolor=blue]{hyperref}
\begin{document}
	
	\title{Minimal Dimensions of Maximal Commutative Matrix Algebras and Sharp Courter-like Bounds}
	\author{Ma\l gorzata Nowak-Kępczyk\footnote{Corresponding Author}\\
		$\rule{0pt}{13pt}$Institute of Informatics\\John Paul II Catholic University of Lublin, Poland\\email: \texttt{malnow@kul.pl}}
\date{}	
\maketitle




\begin{abstract}
	Let \(K\) be an algebraically closed field and let \(M_n(K)\) denote the algebra of \(n\times n\) matrices over \(K\). A classical problem asks for the minimal possible dimension of a maximal commutative subalgebra \(A\subseteq M_n(K)\).
	
	We determine explicit lower bounds for local maximal commutative subalgebras of \(M_n(K)\) for every nilpotency degree of the radical, refining the classical estimate of Laffey in the range relevant to Courter's phenomenon.
	
	In particular, we prove that
	\(\dim A\ge n\)
	for every maximal commutative subalgebra of \(M_n(K)\) with \(n\le13\).
	
	We show that Courter's classical example of dimension \(13\) in
	\(M_{14}(K)\) is the first possible exceptional case and attains the corresponding local bound.
	
	Finally, we introduce a stack construction which, for every \(n\ge14\),
	explicitly constructs a local maximal commutative subalgebra of
	\(M_n(K)\) of dimension strictly smaller than \(n\).
\end{abstract}






\paragraph{Keywords}
maximal commutative subalgebra; matrix algebra; minimal dimension; centralizer; local algebra; nilpotent algebra

\paragraph{MSC}
15A30, 15A27, 16S50	


\section{Introduction}

Let $K$ be an algebraically closed field and let $M_n(K)$ denote the algebra of all $n\times n$ matrices over $K$. A commutative subalgebra $A\subseteq M_n(K)$ is called \emph{maximal commutative} if it coincides with its centralizer in $M_n(K)$, equivalently if $C(A)=A$.

Maximal commutative subalgebras of matrix algebras form a classical topic in linear algebra and associative algebra. Large-dimensional examples are well understood \cite{Brown1997}, while the opposite problem — the existence of maximal commutative algebras of unexpectedly small dimension — is considerably more subtle.

The algebra of diagonal matrices gives the basic example of a maximal commutative subalgebra of dimension $n$. In 1961, Courter \cite{Courter1961} constructed a maximal commutative subalgebra $A\subseteq M_{14}(K)$ with $\dim A=13<n$, showing for the first time that maximal commutative algebras may have dimension strictly smaller than the matrix size. No example in smaller dimension is known.


For $n\le 5$, maximal commutative subalgebras over an algebraically closed field are well understood  \cite{Song1999}. In contrast, for $7\le n\le 13$, neither a classification nor a general exclusion of the inequality $\dim A<n$ was previously available.

A general lower bound for dimension was proved by Laffey \cite{Laffey1983} in 1983, namely
\[\dim A>(2n)^{2/3}-1,\]
but this estimate still allows substantially smaller examples than Courter's.

In our earlier paper \cite{NowakKepczyk}, we showed that no maximal commutative subalgebras of dimension below $n$ exist in $M_6(K)$.

The aim of this paper is to understand the local structure of small
maximal commutative subalgebras and to explain the Courter phenomenon
from this perspective.

Since every commutative algebra decomposes as a direct sum of local algebras, it is natural to focus on the local case. Using the radical filtration
of the natural module \(K^n\), we obtain a canonical lower block
triangular realization for local maximal commutative algebras.

This leads to explicit local lower bounds depending on the nilpotency
degree of the radical. 

The case of radical nilpotency degree $2$ is straightforward. If $A \subseteq M_n(K)$ is a local maximal commutative algebra with $J^2=0$, where $J=\operatorname{rad}(A)$, then every product of two nilpotent elements vanishes, so maximality forces $A$ to contain the full rectangular socle block. Consequently,
\[
\dim A = 1+ab,
\qquad a+b=n,
\]
for suitable positive integers $a,b$. Since
\[
ab \ge n-1,
\]
we obtain
\[
\dim A \ge n.
\]
Thus any Courter-like phenomenon can only occur for radical nilpotency degree at least $3$.

We show that nilpotency degree \(3\) is the most
favorable case for the existence of small maximal commutative algebras.

In particular, local maximal commutative algebras with Loewy signature
\[(n_1,n_2,n_3)\]
satisfy strong restrictions, and if
\[n_1=1\quad\text{or}\quad n_3=1,\]
then necessarily
\[\dim A\ge n.\]

Our bounds show that no Courter-like algebra can occur for \(n\le 13\);
hence Courter's classical example in \(M_{14}(K)\) is the first possible
exceptional case in general.

 We introduce a stack construction which, starting from two explicit brick algebras, yields an infinite family of local maximal commutative algebras. In particular, for every $n \ge 14$, this produces a Courter-like algebra in $M_n(K)$.

Thus Courter's example is not an isolated phenomenon, but the first
member of an infinite family.


\subsection*{Courter-like algebras}


\begin{definition}
	A commutative subalgebra
	$\mathcal A\subseteq M_n(K)$
	with unity is called a \emph{maximal commutative subalgebra} if
	\[	C(\mathcal A)=\mathcal A,	\]
	where
	\[	C(\mathcal A)=\{X\in M_n(K): XA=AX \text{ for all } A\in\mathcal A\}\]
	is the centralizer of $\mathcal A$.
\end{definition}

Equivalently, $\mathcal A$ is maximal commutative if it is not properly contained in any larger commutative subalgebra of $M_n(K)$. This standard characterization will be used throughout the paper.\medskip


\begin{definition}
	Let $\mathcal A\subseteq M_n(K)	$ 	be a maximal commutative algebra with unity.
	We say that \(\mathcal A\) is \emph{local} if every element of
	\(\mathcal A\) is either invertible or nilpotent.
	
	Equivalently, \(\mathcal A\) is local if it possesses a unique maximal
	ideal, namely its Jacobson radical
	\[	\operatorname{rad}(\mathcal A),	\]
	and
	\[	\mathcal A/\operatorname{rad}(\mathcal A)\cong K.
	\]\end{definition}

\begin{remark}
	If 	$A\subseteq M_n(K)$
	is a local commutative algebra with unity, then
	\[
	A=KI_n\oplus \operatorname{rad}(A)
	\]
	as a direct sum of vector spaces.
	
	Indeed, since
	\[
	A/\operatorname{rad}(A)\cong K,
	\]
	the semisimple quotient is generated by the identity matrix.
	Thus every element of \(A\) admits a unique decomposition
	\[
	\lambda I_n+R,	\qquad	\lambda\in K,	\quad	R\in\operatorname{rad}(A).
	\]
\end{remark}

\begin{definition}[Loewy signature]
	Let $A\subseteq M_n(K)$ be a local algebra, and let
$J=\operatorname{rad}(A)$ be nilpotent of degree $r$, i.e.
	\[
	J^{r}K^n=0,	\qquad	J^{r-1}K^n\neq0.	\]	
	This gives the Loewy filtration
	\[
	K^n	\supset	JK^n\supset	J^2K^n	\supset	\cdots	\supset	J^{r-1}K^n	\supset	0.
	\]	
	Choosing complements
	\[
	J^{i-1}K^n	=	V_i\oplus J^iK^n,	\qquad	1\le i\le r,
	\]
	we obtain a vector space decomposition
	\[
	K^n	= 	V_1\oplus V_2\oplus\cdots\oplus V_r.
	\]
	
	The \emph{Loewy signature} of \(A\) is
	\[
	\operatorname{sig}_L(A)	=	(n_1,\dots,n_r),
	\]
	where
	\[
	n_i=\dim V_i	=	\dim_K\!\left(J^{i-1}K^n/J^iK^n\right).
	\]
\end{definition}\medskip

\begin{lemma}[Uniqueness of the Loewy signature]
	The Loewy signature 
	\[
	(n_1,\dots,n_r)
	\]
	is uniquely determined by the algebra
	$A\subseteq M_n(K).$
	In particular, it is independent of the choice of basis.
\end{lemma}

\begin{proof}
	The radical $J=\operatorname{rad}(A)$ is intrinsic to \(A\), hence so are the subspaces
	\[
	K^n\supset JK^n\supset\cdots\supset J^rK^n=0.
	\]
	Therefore the dimensions of the successive quotients
	\[
	J^{i-1}K^n/J^iK^n
	\]
	are uniquely determined.
\end{proof}

\begin{definition}[Loewy form]
	Let $A\subseteq M_n(K)$
	be a local commutative algebra with Loewy signature
	\[
	(n_1,\dots,n_r).
	\]
	
	Choosing a basis adapted to the filtration
	\[
	K^n\supset JK^n\supset\cdots\supset J^rK^n=0,
	\]
	the algebra \(A\) is represented by lower block triangular matrices with
	diagonal block sizes
	\[
	n_1,\dots,n_r.
	\]
	
	We call this realization the \emph{Loewy form} of \(A\).
\end{definition}\medskip
%%%%%%%%%%%%%%%%%%%%%%%%%%%%%%%%%%%%
\begin{proposition}[Loewy form structure]
	Let \(A\subseteq M_n(K)\) be a local maximal commutative algebra with Loewy signature
	$(n_1,\dots,n_r).$
	Then in Loewy form,
	\begin{equation}\label{Loewyeq}
	A=KI+\Lambda+\Sigma,
	\end{equation}
	where
	\[	\Sigma=M_{n_r\times n_1}(K)	\]
	is the full socle block, and
	\[	\Lambda=\langle \lambda_1,\dots,\lambda_d\rangle	\]
	is generated by strictly lower block triangular matrices, linearly independent modulo \(\Sigma\). In particular,
	\[
	\dim A=1+d+n_1n_r.
	\]
\end{proposition}

\begin{proof}
	Choose a basis adapted to the Loewy filtration, so that \(A\) is represented by lower block triangular matrices with diagonal block sizes
	$n_1,\dots,n_r.$
	Since \(A\) is local, we have
	\[	A=KI+\operatorname{rad}(A)	\]
	as a vector space, and every element of \(\operatorname{rad}(A)\) is strictly lower block triangular in this form.
	
	Let \(X\) be any matrix supported only in the bottom-left block
	\[
	M_{n_r\times n_1}(K).
	\]
	Then, by block-size considerations, for every strictly lower block triangular element \(R\in \operatorname{rad}(A)\) one has
	\[
	XR=RX=0.
	\]
	Indeed, \(X\) maps the first Loewy layer to the last one, while any radical element maps each layer strictly downward, so no nonzero block product can occur. Hence \(X\) commutes with every element of \(\operatorname{rad}(A)\). It also commutes with scalar matrices, and therefore
	\[
	X\in C(A).
	\]
	Since \(A\) is maximal commutative, \(C(A)=A\). Thus the whole bottom-left block belongs to \(A\):
	\[
	M_{n_r\times n_1}(K)\subseteq A.
	\]
	We denote this block by
	\[
	\Sigma=M_{n_r\times n_1}(K).
	\]
	
	It remains only to choose a complementary subspace to \(\Sigma\) inside \(\operatorname{rad}(A)\). Let
	\[
	\Lambda=\langle \lambda_1,\dots,\lambda_d\rangle
	\]
	be such a complement. Then the \(\lambda_i\) are linearly independent modulo \(\Sigma\), and we obtain the vector space decomposition
	\[
	A=KI+\Lambda+\Sigma.
	\]
	Consequently, we have (\ref{Loewyeq}). Moreover, \(A\) has the lower block triangular form
	\[
	A
	\subseteq
	\begin{pmatrix}
		K & 0 & 0 & \cdots & 0\\
		* & K & 0 & \cdots & 0\\
		* & * & K & \cdots & 0\\
		\vdots & \vdots & \vdots & \ddots & \vdots\\
		M_{n_r\times n_1}(K) & * & * & \cdots & K
	\end{pmatrix}. 
	\]
\end{proof}



%%%%%%%%%%%%%%%%%%%%%%%%%%%%%%%%%%




\begin{remark}
	If
	\[
	A=A_1\oplus\cdots\oplus A_s
	\subseteq
	M_{m_1}(K)\oplus\cdots\oplus M_{m_s}(K)\subseteq M_n(K),
	\]
	then
	\[
	C(A)=C(A_1)\oplus\cdots\oplus C(A_s).
	\]
	Consequently, \(A\) is maximal commutative in \(M_n(K)\) if and only if each
	\(A_i\) is maximal commutative in \(M_{m_i}(K)\).
\end{remark}


\begin{definition}
	A \emph{Courter-like algebra} is a unital maximal commutative subalgebra
	\[A\subseteq M_n(K)	\]
	such that
	\[\dim A<n.\]
\end{definition}\medskip


\subsection{Maximal commutative algebras with radical nilpotent of degree \(3\)}

Let \(A\subseteq M_n(K)\) be a local maximal commutative algebra whose
radical is nilpotent of degree \(3\), with Loewy signature
\((n_1,n_2,n_3)\).

By the Loewy form (\ref{Loewyeq}), we may write
\(A=KI_n+\Lambda+\Sigma\), where \(\Sigma=M_{n_3\times n_1}(K)\), and
the radical generators occur only in the two off-diagonal blocks
\(\Lambda_{21}\) and \(\Lambda_{32}\).

Thus
\begin{equation}\label{Supru}
	A=
	KI_n\oplus
	\left\{
	\begin{pmatrix}
		0&0&0\\
		X_{21}&0&0\\
		X_{31}&X_{32}&0
	\end{pmatrix}
	:
	X_{21}\in\Lambda_{21},\;
	X_{31}\in\Sigma,\;
	X_{32}\in\Lambda_{32}
	\right\}.
\end{equation}
where
\[\Lambda= \Lambda_1\cup \Lambda_2=\langle \lambda_1,\dots,\lambda_d\rangle\]
denote the linear span of the nilpotent generators  occurring in the
blocks \(\Lambda_{21}\) and \(\Lambda_{32}\), 
We have
\begin{equation} \label{dim} \dim A=1+d+n_1n_3.\end{equation}




\begin{proposition}[Centralizer estimate]\label{centralizer_estimate_nil3}
	Let \(A\subseteq M_n(K)\) be a local maximal commutative algebra with radical
	nilpotent of degree \(3\) and Loewy signature \((n_1,n_2,n_3)\). Write
	\(A=KI+\Lambda+\Sigma\), where \(\Sigma=M_{n_3\times n_1}(K)\) and
	\(\dim\Lambda=d\). Then
	\[
	\dim C(A)\ge
	1+n_1n_3+n_2(n_1+n_3)-d n_1n_3.
	\]
\end{proposition}

\begin{proof}
	Consider the vector space \(\mathcal V\) of matrices of the form
	\[
	X=\alpha I+
	\begin{pmatrix}
		0&0&0\\
		U&0&0\\
		W&V&0
	\end{pmatrix},
	\]
	where \(U\in M_{n_2\times n_1}(K)\), \(W\in M_{n_3\times n_1}(K)\),
	\(V\in M_{n_3\times n_2}(K)\), and \(\alpha\in K\). Thus
	\[
	\dim \mathcal V=1+n_1n_3+n_2(n_1+n_3).
	\]
	
	Every element of \(\mathcal V\) commutes with \(KI\) and with the socle block
	\(\Sigma\). Hence, in order for \(X\in \mathcal V\) to commute with \(A\), it is enough
	to impose commutation with a basis
	\(\lambda_1,\dots,\lambda_d\) of \(\Lambda\).
	
	For each generator \(\lambda_i\), define a linear map
	\(T_i:\mathcal V\to M_{n_3\times n_1}(K)\) by \(T_i(X)=[X,\lambda_i]\). This is
	well-defined because the commutator has support only in the bottom-left block.
	Indeed, if
	\[
	\lambda_i=
	\begin{pmatrix}
		0&0&0\\
		A_i&0&0\\
		B_i&C_i&0
	\end{pmatrix},
	\]
	then a direct block multiplication gives
	\[
	[X,\lambda_i]=
	\begin{pmatrix}
		0&0&0\\
		0&0&0\\
		VA_i-C_iU&0&0
	\end{pmatrix}.
	\]
	Therefore \(\operatorname{rank}T_i\le n_1n_3\).
	
	Now
	\[
	C(A)\cap \mathcal V=\bigcap_{i=1}^d \ker T_i.
	\]
	Using the elementary estimate for intersections of kernels, we obtain
	\[
	\dim(C(A)\cap \mathcal V)
	\ge \dim\mathcal  V-\sum_{i=1}^d \operatorname{rank}T_i
	\ge \dim \mathcal V-dn_1n_3.
	\]
	Since \(C(A)\) contains \(C(A)\cap \mathcal V\), the desired lower bound follows.
\end{proof}\medskip


\begin{theorem}[General local bound for nilpotency degree \(3\)]
	Let \(A\subseteq M_n(K)\) be a local maximal commutative algebra with
	radical nilpotent of degree \(3\) and Loewy signature
	\((n_1,n_2,n_3)\).
	
	Then
	\begin{equation}\label{general3}
		\dim A\ge 1+n_1n_3+	\left\lceil	\frac{n_2(n_1+n_3)}{n_1n_3+1}\right\rceil.
	\end{equation}
\end{theorem}

\begin{proof}
	Write \(A=KI_n+\Lambda+\Sigma\), with \(\dim\Lambda=d\).
	
	Since \(A\) is maximal commutative, \(C(A)=A\). By (\ref{dim}) and
	Proposition~\ref{centralizer_estimate_nil3},
	\[	1+d+n_1n_3	\ge	1+n_1n_3+n_2(n_1+n_3)-d\,n_1n_3.
	\]
	
	Hence
	\[
	d(n_1n_3+1)\ge n_2(n_1+n_3),
	\]
	and therefore
	\[
	d
	\ge
	\left\lceil
	\frac{n_2(n_1+n_3)}{n_1n_3+1}
	\right\rceil.
	\]
	
	Substituting into
	\[
	\dim A=1+d+n_1n_3
	\]
	proves (\ref{general3}).
\end{proof}\medskip

%%%%%%%%%%%%%%%%%%%%%%%%%%%%%%%%

%%%%%%%%%%%%%%%%%%%%%%%%%%%%%%%%%%%%


\begin{corollary}\label{n1_or_n3_one}
	Let $A\subseteq M_n(K)$
	be a local maximal commutative algebra with
radical nilpotent of degree 3 and Loewy signature	$(n_1,n_2,n_3).$
	If	$n_1=1$ or $n_3=1,$ then $\dim A\ge n.$
\end{corollary}

\begin{proof}
	Substituting \(n_1=1\) into (\ref{general3}) gives
	\[	\dim A	\ge	1+n_3+	\left\lceil	\frac{n_2(1+n_3)}{n_3+1}	\right\rceil
	\]
and we obtain
	\[	\dim A	\ge	1+n_3+n_2	=	n.\]
	
	The case \(n_3=1\) is analogous. 
\end{proof}\medskip


We see that a local Courter-like algebra with radical nilpotent of degree 3 can occur only when \[	n_1,n_3\ge2.\]
	

\begin{theorem}[Local lower bound for nilpotency degree \(3\)]
	\label{local_nil3_bound}
	Let	$A\subseteq M_n(K)$
	be a local maximal commutative algebra with radical nilpotent of degree \(3\).	Then
	\[
	\dim A\ge
	\begin{cases}
		n, & n<14,\\[2mm]	\displaystyle	\min_{\substack{n_1+n_2+n_3=n\\ n_1,n_3\ge2}}
		\left\{	1+n_1n_3+	\left\lceil	\dfrac{n_2(n_1+n_3)}{n_1n_3+1}	\right\rceil
		\right\}, & n\ge14.
	\end{cases}
	\]
\end{theorem}

\begin{proof}
	By Corollary~\ref{n1_or_n3_one}, if \(n_1=1\) or \(n_3=1\), then
	\(\dim A\ge n\). Hence a Courter-like case can occur only when
	\(n_1,n_3\ge 2\).
	
	For fixed \(n\), the preceding bound \eqref{general3} gives
	\[
	\dim A\ge
	\min_{\substack{n_1+n_2+n_3=n\\ n_1,n_3\ge2}}
	\left\{
	1+n_1n_3+
	\left\lceil
	\dfrac{n_2(n_1+n_3)}{n_1n_3+1}
	\right\rceil
	\right\}.
	\]
	
	For $n\leq 13$ finite check gives $\dim A\geq n$.
\end{proof}
%%%%%%%%%%%%%%%%%%%%%%%%%%%%%%%%%%%%%%%%%%%%%%%%%%%%


\begin{corollary}[The first Courter-like case]
	\label{first_case}
	For small degrees of matrices
	\[	n=6,7,\ldots,13	\]
	the bound gives
	\[	\dim A	\ge	n.	\]	
	Hence no local Courter-like algebra with radical nilpotent of degree \(3\)
	exists in these dimensions.	
	The first value for which the bound allows
	\[	\dim A<n\]	is	for $n=14$ and Loewy signature $(2,10,2)$.
	Dimension 13 is attained by Courter's classical algebra
	\[	\mathcal C\subseteq M_{14}(K).	\]
\end{corollary}



	
	
%%%%%%%%%%%%%%%%%%%%%%%% nil 4
%%%%%%%%%%%%%%%%%%%%%%%%%%%%%%%%%%%%%%%%%%%%%%
\subsection{Maximal commutative algebras with radical nilpotent of degree \(4\)}

\begin{proposition}[Centralizer estimate for nilpotency degree \(4\)]
	\label{centralizer_estimate_nil4}
	Let \(A\subseteq M_n(K)\) be a local algebra with radical nilpotent of
	degree \(4\), Loewy signature \((n_1,n_2,n_3,n_4)\), and decomposition
	\[
	A=KI_n+\Lambda+\Sigma,
	\qquad
	\dim\Lambda=d.
	\]
	
	Then
	\[
	\dim C(A)
	\ge
	1+n_1n_4+(n_2+n_3)(n_1+n_4)+n_2n_3-d\,n_1n_4.
	\]
\end{proposition}

\begin{proof}
	Consider the ambient lower triangular space
	\[
	\mathcal V=
	\left\{
	\alpha I_n+
	\begin{pmatrix}
		0&0&0&0\\
		*&0&0&0\\
		*&*&0&0\\
		*&*&*&0
	\end{pmatrix}
	\right\},
	\]
	adapted to the Loewy decomposition. Its dimension is
	\[
	1+n_1n_4+(n_2+n_3)(n_1+n_4)+n_2n_3.
	\]
	
	By the same argument as in Proposition~\ref{centralizer_estimate_nil3},
	every element of \(\mathcal V\) commutes with \(KI_n\) and with the socle
	block
	\[
	\Sigma=M_{n_4\times n_1}(K).
	\]
	Thus, to lie in \(C(A)\), it is enough to commute with the generators
	\[
	\lambda_1,\dots,\lambda_d
	\]
	of \(\Lambda\).	
	For each generator, the commutator defines a linear map
	\[
	T_i:\mathcal V\to M_{n_4\times n_1}(K),
	\qquad
	T_i(X)=[X,\lambda_i].
	\]
	As before, the commutator is supported only in the socle block, so
	\[
	\operatorname{rank}T_i\le n_1n_4.
	\]
	Hence
	\[
	\dim C(A)\ge \dim\mathcal V-d\,n_1n_4.
	\]
\end{proof}

\begin{corollary}[Maximal commutative case]
	\label{general4}
	Let \(A\subseteq M_n(K)\) be a local maximal commutative algebra with
	radical nilpotent of degree \(4\) and Loewy signature
	\((n_1,n_2,n_3,n_4)\).
	
	Then
	\[
	\dim A
	\ge
	1+n_1n_4+
	\left\lceil
	\frac{(n_2+n_3)(n_1+n_4)+n_2n_3}{n_1n_4+1}
	\right\rceil.
	\]
\end{corollary}

\begin{proof}
	Write \(A=KI_n+\Lambda+\Sigma\), with \(\dim\Lambda=d\).
	
	Since \(C(A)=A\), Proposition~\ref{centralizer_estimate_nil4} and (\ref{dim}) give
	\[
	1+d+n_1n_4
	\ge
	1+n_1n_4+(n_2+n_3)(n_1+n_4)+n_2n_3-d\,n_1n_4.
	\]
	
	Hence
	\[
	d
	\ge
	\left\lceil
	\frac{(n_2+n_3)(n_1+n_4)+n_2n_3}{n_1n_4+1}
	\right\rceil,
	\]
	which yields the claim.
\end{proof}

\begin{corollary}
	Let \(A\subseteq M_n(K)\) be a local maximal commutative algebra with
	radical nilpotent of degree \(4\) and Loewy signature
	\((n_1,n_2,n_3,n_4)\).
	
	If \(n_1=1\) or \(n_4=1\), then
	\[
	\dim A\ge n.
	\]
\end{corollary}

\begin{proof}
	Substituting \(n_1=1\) into (\ref{general4}) gives
	\[
	\dim A
	\ge
	1+n_4+
	\left\lceil
	\frac{(n_2+n_3)(1+n_4)+n_2n_3}{n_4+1}
	\right\rceil
	\ge
	1+n_4+n_2+n_3=n.
	\]
	
	The case \(n_4=1\) is analogous.
\end{proof}

Thus local Courter-like algebras of nilpotency degree \(4\) may occur only
when
\[
n_1,n_4\ge2.
\]

\begin{theorem}[Local lower bound for nilpotency degree \(4\)]
	\label{local_nil4_bound}
	Let \(A\subseteq M_n(K)\) be a local maximal commutative algebra with
	radical nilpotent of degree \(4\). Then
	\[
	\dim A\ge
	\begin{cases}
		n, & n<23,\\[2mm]
		\displaystyle
		\min_{\substack{n_1+n_2+n_3+n_4=n\\ n_1,n_4\ge2}}
		\left\{
		1+n_1n_4+
		\left\lceil
		\dfrac{(n_2+n_3)(n_1+n_4)+n_2n_3}{n_1n_4+1}
		\right\rceil
		\right\},
		& n\ge23.
	\end{cases}
	\]
\end{theorem}

\begin{proof}
	For \(n<23\), this follows by finite computation. The general estimate is
	an immediate consequence of Corollary~\ref{general4}.
\end{proof}
%%%%%%%%%%%%%%%%%%%%%%%%%%%%%%%%%%%%%%%%%%%%%%
\subsection{General local bound and refinement monotonicity}

\begin{theorem}[General local bound]
	\label{local_bound}
	Let \(A\subseteq M_n(K)\) be a local maximal commutative algebra with
	radical nilpotent of degree \(r\ge3\) and Loewy signature
	\((n_1,\dots,n_r)\).
	
	Put
	\[
	S=\sum_{i=2}^{r-1} n_i,	\qquad	P=\sum_{2\le i<j\le r-1} n_i n_j.
	\]
	
	Then
	\[	\dim A	\ge	1+n_1n_r+	\left\lceil	\frac{S(n_1+n_r)+P}{n_1n_r+1}\right\rceil.
	\]
\end{theorem}

\begin{proof}
	By the same argument as in Propositions~\ref{centralizer_estimate_nil3}
	and~\ref{centralizer_estimate_nil4}, consider the ambient lower block
	triangular space adapted to the Loewy decomposition, consisting of scalar
	matrices together with arbitrary strictly lower block triangular matrices
	supported outside the diagonal.
	
	Its dimension is
	\[
	1+n_1n_r+S(n_1+n_r)+P,
	\]
	where
	\[
	S=\sum_{i=2}^{r-1} n_i,	\qquad	P=\sum_{2\le i<j\le r-1} n_in_j.
	\]
	
	Every element of this space commutes with \(KI_n\) and with the socle block
	\[
	\Sigma=M_{n_r\times n_1}(K).
	\]
	Thus only commutation with the \(d=\dim\Lambda\) generators of
	\(\Lambda\) must be imposed.
	
	For each generator, the commutator defines a linear map into the socle
	block, hence contributes at most
	\[
	n_1n_r
	\]
	independent linear conditions. Therefore
	\[
	\dim C(A)\ge 1+n_1n_r+S(n_1+n_r)+P-d\,n_1n_r.
	\]
	
	Since \(A\) is maximal commutative,
	\[	C(A)=A,	\qquad	\dim A=1+d+n_1n_r,
	\]
	so
	\[
	1+d+n_1n_r
	\ge
	1+n_1n_r+S(n_1+n_r)+P-d\,n_1n_r.
	\]
	Equivalently,
	\[
	d(n_1n_r+1)\ge S(n_1+n_r)+P,
	\]
	which yields the claim.
\end{proof}

\begin{corollary}
	If \(A\subseteq M_n(K)\) is local maximal commutative with radical
	nilpotent of degree \(r\ge3\) and \(n_1=1\) or \(n_r=1\), then
	\[
	\dim A\ge n.
	\]
\end{corollary}

\begin{proof}
	Apply Theorem~\ref{local_bound}.
\end{proof}\medskip





\subsection{Signature-free local bounds and comparison with Laffey's estimate}



Since the Loewy signature is generally not known a priori, we now pass to signature-free local bounds obtained by minimizing over all admissible signatures.

 
\begin{corollary}[Signature-free local bound]
	Let \(A\subseteq M_n(K)\) be a local maximal commutative algebra with
	radical nilpotent of degree \(r\ge3\).
	
	Define
	\begin{equation}\label{D_r}
		D_r(n)
		=
		\min_{\substack{
				n_1+\cdots+n_r=n\\
				n_1,n_r\ge2
		}}
		F_r(n_1,\ldots,n_r),
	\end{equation}
	where \(F_r\) is the lower bound from Theorem~\ref{local_bound}.
	
	Then
	\[
	\dim A\ge D_r(n).
	\]
\end{corollary}

\begin{proof}
	Immediate from Theorem~\ref{local_bound}.
\end{proof}\medskip


\begin{figure}[ht]
	\centering
	\includegraphics[width=0.9\textwidth]{Courters2.png}
	\caption{
		Comparison of the signature-free local lower bounds
		\(D_r(n)\) for nilpotency degrees \(r=3,4,5,6\), together with
		Laffey's estimate and the trivial bound \(\dim A=n\).
		The first dimensions where the bounds allow Courter-like algebras are
		\((14,13)\), \((23,22)\), \((29,28)\), and \((36,35)\),
		respectively.
	}
	\label{Courters}
\end{figure}

Figure~\ref{Courters} suggests the increasing rigidity of local maximal commutative algebras as the nilpotency degree grows: Courter-like examples appear later for higher nilpotency degrees.


\begin{remark}[Persistence of the Courter phenomenon]
	Figure~\ref{Courters} shows that although increasing the nilpotency	degree makes local maximal commutative algebras progressively more	rigid, the Courter phenomenon does not disappear.
	
	The first dimensions where the present signature-free local bounds allow
	\[
	\dim A<n
	\]
	occur at
	\[
	n=14,\ 23,\ 29,\ 36
	\]
	for nilpotency degrees
	\[
	3,\ 4,\ 5,\ 6,
	\]
	respectively. This suggests that the Courter phenomenon is not confined to the 	nilpotency degree \(3\) case, but persists in higher	nilpotency degrees, with the exceptional dimensions shifted upward.
\end{remark}



\begin{remark}[Laffey-type asymptotics]
	For nilpotency degree \(3\), the local bound
	\[
	F_3(n_1,n_2,n_3)
	=
	1+n_1n_3+
	\left\lceil
	\frac{n_2(n_1+n_3)}{n_1n_3+1}
	\right\rceil
	\]
	has the same asymptotic scale as Laffey's classical estimate.
	
	Indeed, ignoring the ceiling function and considering the symmetric
	regime
	\[
	n_1\approx n_3\approx x,
	\qquad
	n_2\approx n-2x,
	\]
	one obtains
	\[
	F_3\approx x^2+\frac{2n}{x},
	\]
	whose minimum occurs at \(x^3\approx n\). Hence
	\[
	F_3\approx 3n^{2/3}.
	\]
	
	The same heuristic applies to every fixed nilpotency degree \(r\ge3\): 	the additional interaction terms contribute only lower-order corrections, 	so
	\[
	D_r(n)\approx 3n^{2/3}.
	\]
	
	Thus the present local bounds have the same exponent \(2/3\) as
	Laffey's estimate
	\[
	(2n)^{2/3}-1,
	\]
	but with a larger leading constant
	\[
	3>2^{2/3}.
	\]
	This is consistent with their stronger performance in the Courter range.
\end{remark}

%%%%%%%%%%%%%%%%%%%%%%%%%%%%%%%%%
\subsection{Global non-local bounds}

We now pass from local algebras to arbitrary maximal commutative
subalgebras of \(M_n(K)\).\medskip


\begin{remark}
	Any signature-free local lower bound immediately globalizes by additivity
	of dimensions under direct sums. More precisely, if
	\[
	A=A_1\oplus\cdots\oplus A_s,
	\]
	where $A_i\subseteq M_{m_i}(K)$
	are local maximal commutative algebras and
	\[
	m_1+\cdots+m_s=n,
	\]
	then
	\[
	\dim A=\sum_{i=1}^s \dim A_i.
	\]	
	Hence any lower bounds for the summands yield a global lower bound for
	\(A\).
\end{remark}

\begin{remark}
	The asymptotic estimate
	\[
	D_3(n)\approx 3n^{2/3}
	\]
	suggests that passing to direct sums should not improve extremal
	dimensions.
	
	Indeed, if
	\[
	z(n)=n-3n^{2/3}
	\]
	models the dimension saving, then
	\[
	z(a+b)-z(a)-z(b)
	=
	3\bigl(a^{2/3}+b^{2/3}-(a+b)^{2/3}\bigr)>0.
	\]
	Thus the asymptotic saving is larger for one local algebra of size
	\(a+b\) than for a direct sum of smaller blocks.
	
This suggests that the global extremal problem may be essentially local.
\end{remark}

\begin{corollary}
	By the local bounds and the direct-sum decomposition, no maximal
	commutative algebra \(A\subseteq M_n(K)\) with \(\dim A<n\) exists for
	\(n\le13\).
\end{corollary}

\begin{proof}
	This follows from the finite verification of the local bounds for
	\(n\le13\), together with additivity of dimension under direct sums.
\end{proof}
	%%%%%%%%%%%%%%%%%%%%%%%%%%%%%%%%%%%%%%
	
%%%%%%%%%%%%%%%%%%%%%%%%%%%%%%%%%
\section{Stack algebras}

The previous sections show that nilpotency degree \(3\) yields the
smallest local lower bounds for Courter-like algebras. The classical
Courter example belongs to this class, with Loewy signature \((2,10,2)\).


 We now introduce a stacking construction for local commutative algebras
of nilpotency degree \(3\). Unlike the ordinary direct sum, stacked
algebras share the same identity and socle, which reduces the resulting
dimension.

\begin{definition}[Stack algebra]
	Let
	\[
	\mathcal A_i\subseteq M_{n_1+s_i+n_3}(K),
	\qquad i=1,2,
	\]
	be local commutative algebras of Loewy signatures
	\[
	(n_1,s_i,n_3),
	\]
	written in Loewy form as
	\[
	\mathcal A_i = KI+\Lambda_i+\Sigma,
	\]
	where
	\[
	\Sigma=M_{n_3\times n_1}(K)
	\]
	is the common socle block.
	
	The stack algebra
	\[
	\mathcal A_1\star\mathcal A_2
	\subseteq
	M_{n_1+s_1+s_2+n_3}(K)
	\]
	is obtained by embedding the generator spaces
	\(\Lambda_1\) and \(\Lambda_2\) into disjoint middle blocks with respect
	to the decomposition
	\[
	K^{n_1}\oplus K^{s_1}\oplus K^{s_2}\oplus K^{n_3}.
	\]
\end{definition}

\begin{proposition}
	The stack algebra
	\[	\mathcal A_1\star\mathcal A_2	\]
	is a local commutative algebra of Loewy signature
	\[	(n_1,s_1+s_2,n_3).	\]
	
	Moreover,
	\[	\dim(\mathcal A_1\star\mathcal A_2)	=
	\dim\mathcal A_1+\dim\mathcal A_2-(1+n_1n_3).	\]
\end{proposition}

\begin{proof}
	Since the generators coming from \(\mathcal A_1\) and \(\mathcal A_2\)
	act on disjoint middle blocks, all mixed products vanish, while the
	original commutativity relations inside each family are preserved.
	Hence
	\[
	\mathcal A_1\star\mathcal A_2
	\]
	is commutative.
	
	Its Loewy signature is
	\[
	(n_1,s_1+s_2,n_3),
	\]
	since the outer layers remain unchanged and the middle layers combine.
	
	Finally, the two algebras share exactly the scalar identity and the
	common socle of dimension \(n_1n_3\), so
	\[
	\dim(\mathcal A_1\star\mathcal A_2)	=
	\dim\mathcal A_1+\dim\mathcal A_2-(1+n_1n_3).
	\]
\end{proof}


\subsection{The first brick: an algebra of signature \((2,5,2)\)}

Let
\[
\mathcal E\subseteq M_9(K)
\]
be the algebra
\[
\mathcal E
=
KI+\Sigma+\langle \lambda_1,\lambda_2,\lambda_3,\lambda_4\rangle,
\]
where
\[
\Sigma=
\left\{
\begin{pmatrix}
	0&0&0\\	0&0&0\\	S&0&0
\end{pmatrix}
:\;
S\in M_2(K)
\right\},
\]
and
\[
\lambda_i=
\begin{pmatrix}
	0&0&0\\	A_i&0&0\\	0&B_i&0
\end{pmatrix},
\]
with respect to the decomposition
\[
K^9=K^2\oplus K^5\oplus K^2.
\]

	Take
\[ A_1= \begin{pmatrix} 1 & -1\\ 0 & 0\\ -1 & 0\\ 1 & 0\\ 0 & -1 \end{pmatrix},\quad   
   A_2= \begin{pmatrix} 0 & 0\\ 0 & 1\\ 0 & 0\\ 0 & 1\\ -1 & 2 \end{pmatrix},\quad  
   A_3= \begin{pmatrix} 1 & 0\\ 0 & 1\\ 0 & 1\\ 0 & -1\\ 1 & 0 \end{pmatrix},\quad  
   A_4= \begin{pmatrix} 0 & 1\\ -1 & -1\\ 0 & -1\\ 0 & 2\\ 0 & 0 \end{pmatrix} \] 
\[ B_1= \begin{pmatrix} 0 & 0 & 1 & 0 & 0\\ 1 & 0 & -1 & -1 & 0 \end{pmatrix} ,\quad  
   B_2= \begin{pmatrix} -1 & 0 & 0 & 1 & 1\\ -1 & 0 & 1 & 2 & 2 \end{pmatrix} \] 
\[ B_3= \begin{pmatrix} -1 & -1 & -1 & 0 & 0\\ 1 & 0 & 1 & 1 & -1 \end{pmatrix} ,\quad  
   B_4= \begin{pmatrix} 1 & 0 & 2 & 1 & 0\\ 0 & 2 & 0 & 0 & 0 \end{pmatrix} \]

Then
\[
B_iA_j=B_jA_i\qquad\text{for all }i,j,
\]
so the generators commute modulo the socle. Hence
$\mathcal{E}$
is commutative.

Moreover,
\[
\dim \langle \lambda_1,\lambda_2,\lambda_3,\lambda_4\rangle=4,
\]
and therefore
\[
\dim \mathcal E=1+4+4=9.
\]

A straightforward symbolic computation shows that
\[
\dim C(\mathcal E)=9.
\]
Hence $C(\mathcal E)=\mathcal E, $
so
$\mathcal{E}$
is maximal commutative. Thus
$\mathcal{E}$
is a local maximal commutative algebra of Loewy signature
$(2,5,2),$ with $d=\dim\Lambda=4=s-1.$

By the local lower bounds established earlier, no algebra of signature
$(2,s,2)$ with $s<5$ can satisfy $d<s.$ Thus $\mathcal{E}$ is the smallest brick of signature \((2,s,2)\) for which the stack construction yields a dimension gain.

%%%%%%%%%%%%%%%%%%%%%%%%%%%%%%%
\subsection{The second brick: a one-dimensional middle brick}

Besides the algebra \(\mathcal E\), we use an elementary brick adding one
middle dimension between the same outer layers \(K^2\) and \(K^2\). 

\[
\mathcal D=
\left\{
\begin{pmatrix}
	a_1&0&0&0&0\\
	0&a_1&0&0&0\\
	a_2&0&a_1&0&0\\
	a_3&a_4&0&a_1&0\\
	a_5&a_6&a_2&0&a_1
\end{pmatrix}
: a_1,\dots,a_6\in K
\right\}
\subseteq M_5(K).
\]

It is immediate that \(\mathcal D\) is commutative, local, and has dimension
\(6\). A direct centralizer computation gives \(C(\mathcal D)=\mathcal D\),
so \(\mathcal D\) is maximal commutative.

It will be used as a one-dimensional middle brick in the stacking construction.


%%%%%%%%%%%%%%%%%%%%%%%%%%%%%%


\begin{lemma}[Mixed rigidity of the bricks]\label{mixed_rigidity}
	Let \(\mathcal P,\mathcal Q\in\{\mathcal E,\mathcal D\}\).
	Then the only linear map
	\[
	T:K^{s_{\mathcal Q}}\to K^{s_{\mathcal P}}
	\]
	satisfying
	\[
	T A_\beta^{\mathcal Q}=0,
	\qquad
	B_\alpha^{\mathcal P}T=0
	\]
	for all generators is \(T=0\).
\end{lemma}

\begin{proof}
	The only point requiring verification is the vanishing of the mixed
	centralizer blocks between distinct stacked components.	For the four possible pairs
	\[
	(\mathcal E,\mathcal E),\quad
	(\mathcal E,\mathcal D),\quad
	(\mathcal D,\mathcal E),\quad
	(\mathcal D,\mathcal D),
	\]
	this reduces to finite homogeneous linear systems for the corresponding
	mixed block \(T\), given by the conditions
	\[
	T A_\beta^{Q}=0,
	\qquad
	B_\alpha^{P}T=0,
	\]
	for all generators of the pair \((P,Q)\).	
	The coefficient matrices of these systems have ranks
	\[
	25,\quad 5,\quad 5,\quad 1,
	\]
	respectively, equal to the numbers of unknowns. Hence each system has only
	the trivial solution \(T=0\), so all mixed centralizer blocks vanish.
	
	The explicit rank computation is given in Appendix~A.
\end{proof}

\begin{theorem}[Maximality of stacked brick algebras]
	For all integers \(q\ge 1\) and \(r\ge 0\), the algebra
	\[
	\mathcal A_{q,r}=\mathcal E^{\star q}\star\mathcal D^{\star r}
	\]
	is local maximal commutative. It has Loewy signature
	\[
	(2,5q+r,2)
	\]
	and dimension
	\[
	\dim \mathcal A_{q,r}=5+4q+r.
	\]
\end{theorem}

\begin{proof}
	By repeated application of the stack construction proposition,
	$\mathcal A_{q,r}$
	is a local commutative algebra with Loewy signature $(2,5q+r,2).$
	
	Since each copy of \(\mathcal E\) contributes four radical generators,
	while each copy of \(\mathcal D\) contributes one, and all stacked
	algebras share the same scalar identity and the same socle of dimension
	\(4\), we obtain
	\[
	\dim \mathcal A_{q,r}=1+4+4q+r=5+4q+r.
	\]
	
	It remains to prove maximal commutativity.
	
	Let
	\[
	X\in C(\mathcal A_{q,r}).
	\]
	With respect to the decomposition
	\[
	K^2
	\oplus
	K^5
	\oplus\cdots\oplus
	K^5
	\oplus
	K
	\oplus\cdots\oplus
	K
	\oplus
	K^2,
	\]
	commutation with the common socle forces \(X\) to preserve the Loewy
	filtration and to have equal scalar action on the outer diagonal blocks.
	
	Now consider the off-diagonal blocks between two distinct brick blocks.
	Such a block
	\[
	T:K^5\to K^5
	\]
	must satisfy
	\[
	TA_i=0,
	\qquad
	B_iT=0,
	\qquad i=1,2,3,4,
	\]
	because it must commute with all generators of the corresponding copy of
	\(\mathcal E\).
	By the rigidity lemma,
	\[
	T=0.
	\]
	Hence no nontrivial mixed blocks between brick components occur.
	
The same argument applies to mixed blocks involving \(\mathcal D\), which vanish
by the corresponding cases of Lemma~23. Therefore the centralizer splits
blockwise into the individual centralizers of the constituent bricks.
	
	Since $	C(\mathcal E)=\mathcal E$
	and $C(\mathcal D)=\mathcal D,$
	it follows that
	\[
	C(\mathcal A_{q,r})=\mathcal A_{q,r}.
	\]
	
	Thus $\mathcal A_{q,r}$
	is maximal commutative.
\end{proof}

%%%%%%%%%%%%%%%%%%%%%%%%%%


%%%%%%%%%%%%%%%%%%%%%%%%%%%%%%%%%%%%%% NO GUARANTEE

\begin{remark}
	It is natural to ask whether the stack construction preserves maximal
	commutativity for arbitrary local maximal commutative algebras of nilpotency
	degree \(3\) with compatible outer layers. We do not use such a general
	statement here. In the construction below, maximality is verified directly for
	the explicit brick algebras by checking the vanishing of all mixed centralizer
	blocks; see Lemma~\ref{mixed_rigidity}.
\end{remark}

%%%%%%%%%%%%%%%%%%%%%%%%%%%%%%%%%%%%%%%%%%%%%%%%%%%%




\section{Construction of Courter-like algebras}

We now show that Courter-like algebras occur in every dimension $n\ge14.$
The construction is based on stacking the brick algebra \(\mathcal E\) and
the diagonal brick \(\mathcal D\).

\begin{theorem}[Construction of Courter-like algebras]
	For every \(n\ge14\), write
	\[	n-4=5q+r,	\qquad	q\ge2,	\quad	0\le r<5.
	\]
	Let \(\mathcal E\) be the brick algebra of signature \((2,5,2)\), and let \(\mathcal D\) denote the brick algebra with one-dimensional middle layer introduced above. Define
	\[	\mathcal A_n	=	\mathcal E^{\star q}\star \mathcal D^{\star r}	\subseteq M_n(K).
	\]
	Then \(\mathcal A_n\) is a local maximal commutative algebra of Loewy
	signature
	\[
	(2,n-4,2)
	\]
	and
	\[	\dim \mathcal A_n	=	5+	\left\lceil	\frac{4(n-4)}5	\right\rceil.
	\]
	In particular, for every \(n\ge14\)
	\[	\dim\mathcal A_n<n.	\]
	
\end{theorem}

\begin{proof}
	By repeated application of the stack construction theorem, $\mathcal A_n$
	is local maximal commutative. The signature is
	\[	(2,5q+r,2)=(2,n-4,2).	\]
	
	Each copy of \(\mathcal E\) contributes \(4\) generators, and each copy
	of \(\mathcal D\) contributes \(1\). Hence
	\[	\dim \mathcal A_n=5+4q+r.	\]
	
	Since \(0\le r<5\),
	\[	\left\lceil	\frac{4(5q+r)}5	\right\rceil	=	4q+r.	\]
	Therefore
	\[	\dim \mathcal A_n=5+\left\lceil	\frac{4(n-4)}5	\right\rceil.	\]
	
	Finally, for \(n\ge14\),
	\[
	5+
	\left\lceil
	\frac{4(n-4)}5
	\right\rceil<n,
	\]
	so every \(\mathcal A_n\) is Courter-like.
\end{proof}

\begin{remark}
	These examples are not claimed to be globally minimal for all large
	\(n\). For larger matrix sizes, smaller dimensions may be obtained by
	using signatures with larger outer Loewy layers. The present construction
	shows instead that the Courter phenomenon persists in every dimension
	\(n\ge14\).
\end{remark}



\begin{table}
	\caption{
		Explicit stack constructions showing that Courter-like algebras exist
		in every dimension \(14\le n\le 28\). The constructions use the brick
		algebra \(\mathcal E\), the diagonal brick \(\mathcal D\).
	}
\[
\begin{array}{c|c|c|c}
	n &  Bound & \mathrm{Optimal\; stack\; construction\;of\;} A_n& \dim A_n \\
	\hline
	14 & 13 & \mathcal E^{\star 2} & 13\\
	15 & 14 & \mathcal E^{\star 2}\star \mathcal D & 14\\
	16 & 15 & \mathcal E^{\star 2}\star \mathcal D^{\star 2} & 15\\
	17 & 16 & \mathcal E^{\star 2}\star \mathcal D^{\star 3} & 16\\
	18 & 17 & \mathcal E^{\star 2}\star \mathcal D^{\star 4} & 17\\
	19 & 17 & \mathcal E^{\star 3} & 17\\
	20 & 18 & \mathcal E^{\star 3}\star \mathcal D & 18\\
	21 & 19 & \mathcal E^{\star 3}\star \mathcal D^{\star 2} & 19\\
	22 & 20 & \mathcal E^{\star 3}\star \mathcal D^{\star 3} & 20\\
	23 & 21 & \mathcal E^{\star 3}\star \mathcal D^{\star 4} & 21\\
	24 & 21 & \mathcal E^{\star 4} & 21\\
	25 & 22 & \mathcal E^{\star 4}\star \mathcal D & 22\\
	26 & 23 & \mathcal E^{\star 4}\star \mathcal D^{\star 2} & 23\\
	27 & 24 & \mathcal E^{\star 4}\star \mathcal D^{\star 3} & 24\\
	28 & 25 & \mathcal E^{\star 4}\star \mathcal D^{\star 4} & 25
\end{array}
\]\end{table}




\begin{remark}
	Although the stack construction produces an explicit infinite family of
	Courter-like algebras, its asymptotic efficiency is not optimal.
	
	Indeed, for the family constructed above we have
	\[
	\dim A_n \sim \frac45\,n
	\qquad (n\to\infty).
	\]
	
	In contrast, the signature-free local bounds suggest a substantially
	smaller asymptotic scale of order $	n^{2/3}.$
	
	Thus the present construction establishes the persistence of the Courter
	phenomenon in infinitely many dimensions, but does not yet capture the
	asymptotically smallest maximal commutative algebras.
\end{remark}

\section*{Conclusion and further directions}

We established explicit local lower bounds for maximal commutative
subalgebras of matrix algebras in terms of the nilpotency degree of the
radical, and showed that nilpotency degree \(3\) yields the most
favorable bounds for the appearance of small-dimensional examples.

Since arbitrary maximal commutative algebras decompose into direct sums
of local ones, and asymptotic considerations suggest that direct sums are
unlikely to improve dimension savings, the extremal problem appears to be
essentially local.

In particular, no maximal commutative subalgebras
\(A\subseteq M_n(K)\) with \(\dim A<n\) exist for \(n\le13\), while
Courter's classical example in \(M_{14}(K)\) is the first exceptional
case.

We also introduced a stack construction showing that this phenomenon is
not isolated: for every \(n\ge14\), we explicitly construct a local
maximal commutative subalgebra of \(M_n(K)\) with dimension strictly
smaller than \(n\).

A natural next step is the classification of local maximal commutative
algebras satisfying \(\dim A<n\), and the determination of asymptotically
optimal constructions.



\section*{Appendix: Rank verification for the stack construction}

The following Mathematica code verifies the rank computations used in
Lemma~\ref{mixed_rigidity}. The four ranks correspond to the pairs
\((\mathcal E,\mathcal E)\), \((\mathcal E,\mathcal D)\),
\((\mathcal D,\mathcal E)\), and \((\mathcal D,\mathcal D)\).

\begin{lstlisting}
	ClearAll["Global`*"];
	
	(* Brick algebra E *)
	
	A1 = {{1, -1}, {0, 0}, {-1, 0}, {1, 0}, {0, -1}};
	A2 = {{0, 0}, {0, 1}, {0, 0}, {0, 1}, {-1, 2}};
	A3 = {{1, 0}, {0, 1}, {0, 1}, {0, -1}, {1, 0}};
	A4 = {{0, 1}, {-1, -1}, {0, -1}, {0, 2}, {0, 0}};
	
	B1 = {{0, 0, 1, 0, 0},
		{1, 0, -1, -1, 0}};
	
	B2 = {{-1, 0, 0, 1, 1},
		{-1, 0, 1, 2, 2}};
	
	B3 = {{-1, -1, -1, 0, 0},
		{1, 0, 1, 1, -1}};
	
	B4 = {{1, 0, 2, 1, 0},
		{0, 2, 0, 0, 0}};
	
	AE = {A1, A2, A3, A4};
	BE = {B1, B2, B3, B4};
	
	(* Diagonal brick D *)
	
	AD = {{{1, 0}}};
	BD = {{{1}, {0}}};
	
	(* Rank of the mixed rigidity system:
	T A_j^Q = 0,
	B_i^P T = 0
	*)
	
	mixedRank[AP_, BP_, AQ_, BQ_] := Module[
	{sP, sQ, vars, T, equations, coeffs},
	
	sP = Length[AP[[1]]];
	sQ = Length[AQ[[1]]];
	
	vars = Array[t, {sP, sQ}];
	T = vars;
	
	equations = Join[
	Flatten[Table[T . AQ[[j]], {j, Length[AQ]}]],
	Flatten[Table[BP[[i]] . T, {i, Length[BP]}]]
	];
	
	coeffs = Normal[
	CoefficientArrays[equations, Flatten[vars]][[2]]
	];
	
	MatrixRank[coeffs]
	];
	
	rankEE = mixedRank[AE, BE, AE, BE];
	rankED = mixedRank[AE, BE, AD, BD];
	rankDE = mixedRank[AD, BD, AE, BE];
	rankDD = mixedRank[AD, BD, AD, BD];
	
	{
		{"pair", "rank", "unknowns"},
		{"(E,E)", rankEE, 25},
		{"(E,D)", rankED, 5},
		{"(D,E)", rankDE, 5},
		{"(D,D)", rankDD, 1}
	}
\end{lstlisting}


\section*{Declarations}

\textbf{Funding} \\
The author received no external funding.

\textbf{Conflicts of interest} \\
The author declares that there is no conflict of interest.

\textbf{Data availability} \\
Data sharing is not applicable to this article as no datasets were generated or analyzed.

\textbf{Author contributions} \\
The author contributed solely to all aspects of this work.


\begin{thebibliography}{9}
	
	
	
	\bibitem{Brown1997}
	W. C. Brown,
	\emph{On maximal commutative subalgebras of matrix algebras},
	J. Algebra \textbf{187} (1997), 706--722.
	doi:10.1006/jabr.1996.6789
	
		\bibitem{Courter1961}
	R. C. Courter,
	\emph{Maximal commutative subalgebras of matrix algebras},
	Duke Math. J. \textbf{28} (1961), 225--232.
	doi:10.1215/S0012-7094-61-02827-9
	
	\bibitem{CurtisReiner}
	C. W. Curtis and I. Reiner,
	\emph{Representation Theory of Finite Groups and Associative Algebras},
	Wiley, New York, 1962.
	
\bibitem{AndersonFuller}
F. W. Anderson and K. R. Fuller,
\emph{Rings and Categories of Modules},
Graduate Texts in Mathematics, Vol. 13,
Springer, New York, 1992.
	
		\bibitem{Song1999}
	Y. Song,
	\emph{Notes on the constructions of maximal commutative subalgebras of $M_n(k)$},
	Comm. Algebra \textbf{27} (1999), no. 8, 4163--4173.
	doi:10.1080/00927879908826670
	
		\bibitem{Laffey1983}
	T. J. Laffey,
	\emph{On commutative subalgebras of matrix algebras},
	Linear Algebra Appl. \textbf{55} (1983), 1--7.
	doi:10.1016/0024-3795(83)90112-3
	
	\bibitem{NowakKepczyk}
	M. Nowak-K\k{e}pczyk,
	\emph{On the minimal dimension of maximal commutative subalgebras of {$M_6(k)$}},
	arXiv:2604.23322 (2026), doi:10.48550/arXiv.2604.23322.
	
	\bibitem{SuprunenkoTyshkevich}
	D. A. Suprunenko, R. I. Tyshkevich,
	\emph{Commutative Matrices},
	Academic Press, New York, 1968.
	
	\bibitem{ArzhantsevZaitseva}
	I. Arzhantsev, Y. Zaitseva, 
	\emph{Equivariant completions of affine spaces}, Russian Math. Surveys \textbf{77} (2022), no. 4, 571--622. doi:10.1070/RM10040
	

	

	
\end{thebibliography}

\end{document}
